% !PRINTED 598
\weberpage{598}
%<<CONT
call a \gesperrt{metacyclic equation}. The cyclic equations
themselves are contained among them as a special case,
and by what was remarked in \S.\,175 the metacyclic
equations are the same as the equations solvable by radicals. If $P$ is the group of such an equation, then by
the preceding section it must possess a normal divisor of prime index $j_1$, which we now wish to denote by $P_1$.
If $P_1$ consists of the single identical permutation, then $P$
is itself cyclic and of prime degree. If $P_1$ is not the unit group, then $P_1$ must in turn contain a normal divisor $P_2$ of prime index $j_2$, and so on, until we finally arrive at the unit
group.

That this composition of its group $P$ is also the sufficient condition for a metacyclic equation follows from the preceding section. We
therefore pronounce the theorem:

\begin{quote}
\noindent III.\ \gesperrt{The necessary and sufficient condition for
a metacyclic equation is that there is
a sequence of groups}

\[
P, P_1, P_2, P_3 \ldots
\]

\gesperrt{whose first is the Galois group of the
equation, whose last is the unit group,
and of which each following one is a normal divisor
of the one next preceding, of prime index.}
\end{quote}

Accordingly we call a permutation group $P$ which has this
property, for which therefore a chain of groups

\[
P, P_1, P_2 \ldots P_{\mu-1}, 1
\]

can be so determined that each term is a normal divisor of the preceding, of prime index, a \gesperrt{metacyclic
group}\footnote{The expression ``metacyclic groups'' was first used by \gesperrt{Kronecker},
though in a more restricted sense. I should like here to propose this
slight generalisation of an already known expression in place of the
expression ``solvable groups'' used by \gesperrt{Frobenius} and \gesperrt{Hölder}.}.

We have here obtained the condition for the \gesperrt{complete} solvability of an equation by a chain of cyclic
equations. But it is still a question of the further inquiry,
%CONT>>
